% Part: incompleteness
% Chapter: representability-in-q
% Section: c-representable

\documentclass[../../../include/open-logic-section]{subfiles}

\begin{document}

\olfileid{inc}{req}{cre}
\olsection{تابع‌های $C$ در $\Th{Q}$ بازنمایی‌پذیرند}

باید نشان دهیم هر تابع در $C$ در $\Th{Q}$ بازنمایی‌پذیر است. در پایان
باید نشان دهیم چگونه به هر تابع $k$موضعی
$f(x_0,\dots,x_{k-1})$ در $C$، فرمول
$!A_f(x_0,\dots,x_{k-1},y)$ را نسبت دهیم که آن را بازنمایی کند.

\begin{lem}
برای اعداد طبیعی $n$ و $m$، اگر $n \neq m$، آنگاه
$\Th{Q} \Proves \num n \neq \num m$.
\end{lem}

\begin{proof}
با استقرا بر $n$ نشان دهید برای هر $m$، اگر $n \neq m$، آنگاه
$\Th{Q} \Proves \eq/[\num n][\num m]$.

در حالت پایه $n = 0$. اگر $m$ برابر $0$ نباشد، برای عدد طبیعی~$k$ای
$m = k + 1$. اصل موضوعی داریم که می‌گوید
$\lforall[x][\eq/[0][x']]$. با یک اصل موضوع سور و با جایگزینی $x$ با
~$\num k$ نتیجه می‌گیریم $\eq/[0][\num k']$. اما $\num k'$ همان
$\num m$ است.

در گام استقرا می‌توانیم صدق ادعا را برای $n$ فرض و $n+1$ را بررسی
کنیم. بگذارید $m$ هر عدد طبیعی باشد. دو حالت ممکن است: یا $m = 0$، یا
برای $k$ای $m = k+1$. در حالت نخست، اصل موضوع جانشین‌نبودن صفر مستقیماً
می‌دهد $\Th{Q} \Proves \eq/[\num n'][0]$. در حالت دوم، فرض کنید
$n+1 \neq k+1$. آنگاه $n \neq k$. بنا بر فرض استقرا برای $n$،
$\Th{Q} \Proves \eq/[\num n][\num k]$. اصلی موضوع داریم که می‌گوید
$\lforall[x][\lforall[y][\eq[x'][y'] \lif \eq[x][y]]]$. با یک اصل موضوع
سور داریم
$\eq[\num n'][\num k'] \lif \eq[\num n][\num k]$. با منطق گزاره‌ای در
$\Th{Q}$ می‌توانیم نتیجه بگیریم
$\eq/[\num n][\num k] \lif \eq/[\num n'][\num k']$. با وضع مقدم
$\eq/[\num n'][\num k']$ به دست می‌آید، که همان نتیجهٔ مطلوب است، زیرا
$\num k'$ برابر $\num m$ است.
\end{proof}

توجه کنید این لم چندان پرمحتوا نیست: در اصل می‌گوید $\Th{Q}$ می‌تواند
اثبات کند اعدادنمای متفاوت بر اشیای متفاوت دلالت می‌کنند. برای نمونه،
$\Th{Q}$ اثبات می‌کند $0'' \neq 0'''$. اما نشان‌دادن این امر در حالت
کلی قدری دقت می‌طلبد. همچنین توجه کنید با آنکه از استقرا استفاده
می‌کنیم، این استقرا \emph{بیرون} از $\Th{Q}$ است.

می‌توانیم صفر، جانشین، جمع، ضرب، تابع مشخصهٔ برابری و تصویرها را
بازنمایی کنیم. در هر مورد فرمول بازنمای مناسب کاملاً سرراست است؛ برای
نمونه، صفر با فرمول
\[
y = 0
\]
بازنمایی می‌شود، جانشین با فرمول
\[
x_0' = y
\]
و جمع با فرمول
\[
x_0 + x_1 = y.
\]
کار اصلی نشان‌دادن این است که $\Th{Q}$ می‌تواند گزاره‌های مربوط را
اثبات کند. برای مثال، اینکه جمع با فرمول بالا بازنمایی می‌شود مستلزم
آن است که نشان دهیم برای هر جفت عدد طبیعی $m$ و $n$، $\Th{Q}$ اثبات
می‌کند
\[
\eq[\num n + \num m][\num {n+m}]
\]
و
\[
\lforall[y][(\eq[(\num n + \num m)][y] \lif \eq[y][\num {n+m}])].
\]

ترکیب چطور؟ فرض کنید $h$ چنین تعریف شده باشد:
\[
h(x_0,\dots,x_{l-1}) = f(g_0(x_0,\dots,x_{l-1}), \dots,
g_{k-1}(x_0,\dots,x_{l-1})).
\]
فرض کنید پیش‌تر فرمول‌های $!A_f, !A_{g_0}, \dots,!A_{g_{k-1}}$ را
یافته‌ایم که به‌ترتیب تابع‌های $f,g_0,\dots,g_{k-1}$ را بازنمایی
می‌کنند. آنگاه می‌توانیم فرمول $!A_h$ را که $h$ را بازنمایی می‌کند
تعریف کنیم؛ به این صورت که $!A_h(x_0,\dots,x_{l-1},y)$ را برابر عبارت
زیر بگیریم:
\begin{multline*}
  \lexists[z_0]\dots\lexists[z_{k-1}][
  !A_{g_0}(x_0,\dots,x_{l-1},z_0) \land \dots \land\\
  !A_{g_{k-1}}(x_0,\dots,x_{l-1},z_{k-1}) \land
  !A_f(z_0,\dots,z_{k-1},y)].
\end{multline*}

سرانجام جست‌وجوی نامحدود را بررسی کنیم. فرض کنید $g(x,\vec z)$ منظم و
در $\Th{Q}$ با فرمول $!A_g(x,\vec z,y)$ بازنمایی‌پذیر باشد. بگذارید
$f$ با $f(\vec z) = \mu x \; g(x,\vec z)$ تعریف شود. می‌خواهیم فرمول
$!A_f(\vec z,y)$ را بیابیم که $f$ را بازنمایی کند. انتخاب طبیعی چنین
است:
\[
!A_f(\vec z,y) \ident !A_g(y,\vec z,0) \land \lforall[w][(w < y
\lif \lnot !A_g(w,\vec z,0))].
\]

با استفاده از چند لم اضافی می‌توان نشان داد این فرمول کار می‌کند؛ برای
نمونه:

\begin{lem}
  برای هر متغیر $x$ و هر عدد طبیعی $n$، $\Th{Q}$ اثبات می‌کند
  $\eq[(x' + \num n)][(x + \num n)']$.
\end{lem}

باز هم شایان ذکر است که این گزاره ضعیف‌تر از آن است که $\Th{Q}$ فرمول
$\lforall[x][\lforall[y][(x' + y) = (x +y)']]$ را اثبات کند، زیرا
گزارهٔ اخیر کاذب است.

\begin{proof}
اثبات، طبق معمول، با استقرا بر $n$ است. در حالت پایه $n = 0$ باید نشان
دهیم $\Th{Q}$ فرمول $\eq[(x'+0)][(x+0)']$ را اثبات می‌کند. اما داریم:
\begin{eqnarray*}
& & \eq[(x' + 0)][x'] \quad \text{از اصل موضوع ۴} \\
& & \eq[(x + 0)][x] \quad \text{از اصل موضوع ۴} \\
& & \eq[(x+0)'][x'] \quad \text{از سطر ۲} \\
& & \eq[(x' + 0)][(x + 0)'] \quad \text{از سطرهای ۱ و ۳}
\end{eqnarray*}
در گام استقرا می‌توانیم فرض کنیم
$\eq[(x' + \num n)][(x + \num n)']$ را در $\Th{Q}$ استنتاج کرده‌ایم.
چون $\num{n+1}$ همان $\num n'$ است، باید نشان دهیم $\Th{Q}$ فرمول
$\eq[(x' + \num n')][(x + \num n')']$ را اثبات می‌کند. زنجیرهٔ
برابری‌های زیر را می‌توان در $\Th{Q}$ استنتاج کرد:
\begin{eqnarray*}
(x' + \num n') & \eq & (x' + \num n)' \quad \text{اصل موضوع ۵} \\
& \eq & (x + \num n')' \quad \text{از فرض استقرا}
\end{eqnarray*}
\end{proof}

\begin{lem}
\begin{enumerate}
\item $\Th{Q}$ اثبات می‌کند $\lnot (x < \num 0)$.
\item برای هر عدد طبیعی $n$، $\Th{Q}$ اثبات می‌کند
\[
x < \num {n+1} \lif (\eq[x][\num 0] \lor \dots \lor \eq[x][\num n]).
\]
\end{enumerate}
\end{lem}

\begin{proof}
بخش ۱ و قسمتی از بخش ۲ را به‌طور غیرصوری انجام می‌دهیم؛ یعنی فقط
راهنمایی‌هایی برای ساخت !!{derivation} صوری می‌آوریم.

برای بخش ۱، بنا بر تعریف $<$ باید در $\Th{Q}$ فرمول
$\lnot \lexists[y][\eq[(y'+x)][0]]$ را اثبات کنیم؛ این فرمول، با
استفاده از اصول موضوعه و قواعد منطق مرتبهٔ اول، با
$\lforall[y][\eq/[(y' + x)][0]]$ هم‌ارز است. ایده چنین است: فرض کنید
$\eq[(y'+x)][0]$. اگر $x$ برابر $0$ باشد، داریم
$\eq[(y' + 0)][0]$. اما بنا بر اصل موضوع ۴ از $\Th{Q}$،
$\eq[(y' + 0)][y']$، و بنا بر اصل موضوع ۲،
$\eq/[y'][0]$؛ که تناقض است. پس
$\lforall[y][\eq/[(y' +x)][0]]$. اگر $x$ برابر $0$ نباشد، بنا بر اصل
موضوع ۳، $z$ای هست که $\eq[x][z']$. اما آنگاه
$\eq[(y' + z')][0]$. بنا بر اصل موضوع~۵،
$\eq[(y' + z)'][0]$، که باز با اصل موضوع~۲ تناقض دارد.

برای بخش ۲ بر $n$ استقرا کنید. حالت پایه، یعنی $n =0$، را در نظر
بگیریم. در این حالت باید نشان دهیم
$x < \num 1 \lif \eq[x][\num 0]$. فرض کنید $x < \num 1$. آنگاه بنا بر
اصل موضوع معرف $<$، داریم
$\lexists[y][\eq[(y'+x)][0']]$. فرض کنید $y$ این خاصیت را دارد؛ پس
$\eq[y'+x][0']$.

باید نشان دهیم $\eq[x][0]$. بنا بر اصل موضوع ۳، اگر $x$ برابر $0$
نباشد، برای $z$ای برابر $z'$ است. آنگاه
$\eq[(y' + z')][0']$. بنا بر اصل موضوع~۵ از $\Th{Q}$،
$\eq[(y'+z)'][0']$. بنا بر اصل موضوع ۱، داریم
$\eq[(y' + z)][0]$. اما بنا بر تعریف این یعنی $z < 0$، که با بخش~۱
تناقض دارد.
\end{proof}

\begin{explain}
نشان داده‌ایم مجموعهٔ تابع‌های محاسبه‌پذیر را می‌توان به‌صورت مجموعهٔ
تابع‌های بازنمایی‌پذیر در $\Th{Q}$ مشخص کرد. در واقع اثبات کلی‌تر است.
از تعریف بازنمایی‌پذیری به‌آسانی می‌توان دید هر نظریه‌ای که $\Th{Q}$
را گسترش می‌دهد، یا بتوان $\Th{Q}$ را در آن تفسیر کرد، قادر به
بازنمایی تابع‌های محاسبه‌پذیر است. برعکس، در هر دستگاه
!!{derivation}ی که مفهوم اثبات در آن محاسبه‌پذیر باشد، هر تابع
بازنمایی‌پذیر محاسبه‌پذیر است. بنابراین، برای نمونه، مجموعهٔ تابع‌های
محاسبه‌پذیر را می‌توان مجموعهٔ تابع‌های بازنمایی‌شده در حساب پئانو یا
حتی نظریهٔ مجموعه‌های زرملو--فرانکل دانست. چنان‌که گودل اشاره کرد، این
نتیجه تا اندازه‌ای شگفت‌انگیز است. خواهیم دید در بحث اثبات‌پذیری،
پرسش‌ها به نظریه‌ای که در نظر می‌گیریم بسیار حساس‌اند؛ به‌طور تقریبی،
هرچه اصول موضوعه قوی‌تر باشند، چیزهای بیشتری می‌توان اثبات کرد. اما در
گسترهٔ وسیعی از نظریه‌های اصل‌موضوعی، تابع‌های بازنمایی‌پذیر دقیقاً
همان تابع‌های محاسبه‌پذیرند.
\end{explain}

\end{document}
